\section{The Irish Capuchin Network}

The first resident phase belonged to a religious-order system wider than Willits. Irish Capuchins staffed multiple Mendocino County centers, exchanged assignments, traveled between communities, and reported to both their superiors and the archbishop. Their archive preserves that network more clearly than it preserves the daily life of any one congregation.

Ambrose Brunton's appointment in 1923 gave Saint Anthony its resident center. Raphael Quinn remained important from Ukiah: his authority record credits him with the Willits church project as well as the better-dated new Saint Mary's church. Finbarr O'Callaghan served at Willits during the later 1920s. Sebastian Brennan, whose Mendocino ministry had begun elsewhere in 1903, was posted to Willits in October 1928. Clergy movement makes a complete pastor list difficult to reconstruct from isolated annual editions; it also reveals how the order treated county ministry as a shared field rather than a set of sealed parishes.

Photograph descriptions provide rare glimpses of the network occupying a place. A 1926 catalog item shows Archbishop Hanna and Capuchin friars at Saint Anthony. It establishes their presence together but not why they gathered. Nothing in the accessible description warrants calling the occasion a dedication, confirmation, visitation, or anniversary. The absence of an event label is a reason to preserve the image record, not to supply a picturesque scene from imagination.

The most tantalizing Willits records date to 1929. The archive's Willits subseries, IE~CA~WA/7/14, identifies a Golden Jubilee record for Brennan and a ``Parish Status and Financial Statement'' for Saint Anthony. Their existence proves that the parish was reporting institutional and financial condition during Brennan's Willits assignment. Their contents were not accessible to this study. No income, debt, membership, property, or building claim has been taken from the titles alone.

The Capuchin system gave a small town more than a succession of pastors. It linked Willits to Ukiah, the coast, San Francisco ecclesiastical government, the order's Irish province, and rural Catholic communities far beyond the town. It also created an archival bias. Personnel files, superior's correspondence, group photographs, and institutional reports survive because the order generated and preserved them. The cooking, cleaning, fundraising, transport, lodging, catechesis, music, and informal care supplied by laypeople---especially women whose names seldom enter directories---are less visible.

That imbalance prevents a clerical roster from becoming the history itself. Priests made sacramental travel possible and their assignments mark changes in governance. A parish survived between visits, appointments, and annual reports because local people maintained a place and expected ministry to return. The 1923 deputation is unusually valuable precisely because it lets those local demands break through an archive otherwise organized around clergy.

\evidenceframe{Irish Capuchin Archives authority records for Ambrose Brunton, Raphael Quinn, Sebastian Brennan, and other friars; photographic item descriptions; Willits subseries IE~CA~WA/7/14; annual Catholic directories.}{The resident parish formed within a countywide Irish Capuchin personnel and mission network, with identifiable assignments, church responsibility, episcopal contact, and institutional reporting.}{Archive descriptions are not the full documents. They do not establish complete pastor succession, the purpose of the 1926 gathering, the contents of the 1929 reports, or lay experience.}
